\chapter{Supersonic Flow around a sphere}

Experimental data for compressible flow around a sphere for different Mach
numbers and Reynolds numbers are available over the number of publications
starting from the 18th century. A review article by Miller \&
Bailey \cite{Miller:1979} studies the experiments of the 18th and 19th
centuries for different Mach numbers $0.2
\leq M \leq 2.0$ at $Re \approx 10^7$. The analyses in the mentioned article
show that the experimental data obtained three centuries ago are in good
agreement to the modern data. 

Sphere measurements obtained by \cite{Miller:1979} and \cite{Bailey:1971}
show that regardless of the diameter size of the sphere, the drag coefficient
increases rapidly in the transonic region and for $M\geq1.6$ and $Re\geq 10^5$
it stays to be constant with further increases in Reynolds number. 
Experimental analyses in \cite{Bailey:1971} and earlier in \cite{Charters:1945}
show that the drag coefficient 
slowly decreases when the Mach number increases above $M\geq 2$ for the
flow with $Re \gtrapprox 10^6$. 
We leave a
complete discussion for different Mach numbers for future research, but in this
paper we present results for only one Mach number, in order
to see if the RNS solution gives a correct result according to the experiment.

We use Algorithm \ref{adapt:alg} for the supersonic flow
around a sphere  with diameter $d = 0.074$ at Mach number $M=2$. For this
Mach number the flow characterized by a detached three
dimensional bow shock wave in front of the sphere and is in a mixed
subsonic and supersonic flow behind the sonic line $M=1$. Here the pressure drag
at the stagnation point is high compared to the
subsonic
case. An attached shock
wave develops in the rear of the sphere, which is also counted as a
substantial source of the pressure drag. Therefore, the area for the
pressure stagnation point and attached shock waves should be resolved by
computation for the correct drag coefficient.
We observe from the results that the adaptive RNS method tries to resolve these
regions.

We present the results from the simulation in Table \ref{table:sphere}.
$10\%$ of the largest cells are refined during the adaptive algorithm.
We start with an initial coarse mesh, which has 9 720 vertices and 53 312 cells,
which reaches 168 820 vertices and 918 605 cells after 9th adaptive
iteration. The error bound of the drag coefficient
$\displaystyle\sum_{n,K}\eta_n^K$ converge by the
mesh refinement, as the drag coefficient approximately
stays around the experimental data $C_{dp} \approx 1$. In Figure
\ref{fig:sphere_cd} the drag coefficients are plotted after each adaptive
iterations.  

We plot the solution of the adaptive algorithm in Figure \ref{fig:sphere_sol}
after
eight refinements. The right column of the figure shows the primal solutions and
the left column presents the corresponding dual solutions. The magnitude of the
dual solution is high in the areas with the significant contribution to the
pressure drag force. The plot of the primal pressure shows that already in
eight adaptive step the pressure structures are close to be resolved.

In Figure \ref{fig:mesh_3D} we present the initial coarse mesh and the mesh
after nine adaptive refinements. After each refinement the boundary nodes,
which appear from the Rivara algorithm are projected to the surface
of the sphere. Also, we notice that similar to the 2D result, the algorithm
focuses
to resolve the pressure stagnation point, area of the sonic line and
attached shock wave, which is expected since they are the main source of
pressure drag. With ad hoc refinement, for instance residual
based or gradient based adaptation, the region with strong shocks are well
resolved. However, we notice that the duality based adaptive algorithm does not
resolve a propagating bow shock and other strong discontinuities, it only
refines the areas with the largest error contribution. Consequently,	 it
significantly decreases the computational cost of the drag force computation.

\begin{table}[ht]
\caption{The convergence history of the drag coefficient. 
Here $\bar C_{dp}$ is a mean value of $C_{dp}$ over the time interval
$[t-\epsilon_t, t]$, where $\epsilon_t$ is a small number, and
$\displaystyle\sum_{n, K}\eta_n^K$ denotes a sum of error indicators.}
\centering
\begin{tabular}{ c | c | c | c | c | c  }
\hline
\#iter & \#vertices & \#cells & $S$ &   $\bar C_{dp}$ &
$\displaystyle\sum_{n,K}\eta_n^K$ \\
  \hline 
0 &9720&  53312 &0.8218        &0.8115          &1.5018 \\
1&17405&  94572 &0.9571        &0.8360          &1.3096 \\
2&22915&  124884 &1.1586         &0.8502         &1.0234 \\
3&29686&  160313 &1.4233         &0.9085         &0.8651 \\
4&39328&  212637 &1.6683         &0.9474         &0.7352 \\
5&51746&  280031 &1.8440         &0.9705         &0.6236 \\
6&69418&  376144 &2.0050         &0.9873         &0.5429 \\
7&93696&  508209 &2.0538         &0.9982         &0.4532 \\
8&  126012&  685079 &2.1038      &1.0046         &0.3844 \\
9& 168820 & 918605 & 2.0541      &1.0081         & 0.3199 \\
\hline 
 \end{tabular}\label{table:sphere}
\end{table}
 
 
 \begin{figure}[bhpt]
           \centerline{
           \includegraphics[width=0.45\textwidth]{jpg/sphere/dual_rho.eps}
           \includegraphics[width=0.45\textwidth]{jpg/sphere/density.eps}
           }
	   \vspace{0.1in}
           \centerline{
           \includegraphics[width=0.45\textwidth]{jpg/sphere/dual_pre.eps}
           \includegraphics[width=0.45\textwidth]{jpg/sphere/pressure.eps}
           }
	   \vspace{0.1in}
           \centerline{
           \includegraphics[width=0.45\textwidth]{jpg/sphere/dual_u.eps}
           \includegraphics[width=0.45\textwidth]{jpg/sphere/velocity.eps}
           }
	   \vspace{0.1in}
  \caption{Supersonic flow around a sphere: the dual solution at time $t=0$
of density \textit{(top-left)}, pressure \textit{(middle-left)} and 
magnitude of velocity
\textit{(bottom-left)}; the primal solution at time $t=0.5$ of density
\textit{(top-right)}, pressure \textit{(middle-right)} and 
magnitude of velocity
\textit{(bottom-right)}. The contours are plotted in the collormap. The sonic
line in red is plotted together with the dual solution.
}
  \label{fig:sphere_sol}
\end{figure}

 
 \begin{figure}[bhpt]
           \centerline{
           \includegraphics[width=0.45\textwidth]{jpg/sphere/mesh_it_0.eps}
           \includegraphics[width=0.45\textwidth]{jpg/sphere/mesh_it_4.eps}
           }
	   \vspace{0.1in}
           \centerline{
           \includegraphics[width=0.45\textwidth]{jpg/sphere/mesh_it_9.eps}
           \includegraphics[width=0.45\textwidth]{jpg/sphere/mesh_it_0_yz.eps}
           }
	   \vspace{0.1in}
           \centerline{
           \includegraphics[width=0.45\textwidth]{jpg/sphere/mesh_it_9_yz.eps} 
      \includegraphics[width=0.45\textwidth]{jpg/sphere/mesh_it_9_yz_back.eps}
           }
	   \vspace{0.1in}
  \caption{Supersonic flow around a sphere: the  $(x,y)$ -
view of the mesh for the initial mesh \textit{(top-left)}, four times
\textit{(top-right)}
and nine times \textit{(middle-left)} adaptive refinement according to the
drag
force together, the sonic line is plotted in red, $(y,z)$ - view
of the mesh for $x=x_s$ of the initial mesh \textit{(middle-right)}, where
$x_s$ - is
$x$ coordinate of the center of sphere,  
the finest mesh close to the stagnation point of pressure $x =
x_s - d/2$ \textit{(below-left)}, and the finest mesh at $x =
x_s + 0.02$  from back \textit{(below-right)}
}
  \label{fig:mesh_3D}
\end{figure}

\begin{figure}[bhpt]
           \centerline{          
\includegraphics[width=0.7\textwidth]{jpg/sphere/drag_sphere_M_2_0.eps}
           }
 	\centerline{       
\includegraphics[width=0.7\textwidth]{jpg/sphere/drag_zoom.eps} 
           }
  \caption{
  Supersonic flow around a sphere: the drag coefficient $C_{dp}$ in different
adaptive iterations.}
  \label{fig:sphere_cd}
\end{figure}


